\newpage % Rozdziały zaczynamy od nowej strony.
\section{Wstęp}
\subsection{Sformowanie problemu}
  W ostatnich latach można zaobserwować dynamiczny rozwój technologii sztucznej inteligencji. Technologia ta znajduję coraz szersze zastosowanie w  różnych gałęziach przemysłu, od branży IT aż po medycynę. Narzędzia te wykorzystywane są do automatyzacji procesów oraz wspierania człowieka w wykonywanej pracy. Przemysł rolniczy nie jest w tej kwestii wyjątkiem - także w tej gałęzi gospodarki wprowadzane są nowoczesne technologie, nie będące tylko nowym modelem traktora, lecz szerokim zakresem produktów opartych na uczeniu maszynowym. 
  
  Choroby oraz szkodniki to jedne z głównych przyczyn śmierci roślin uprawnych - niszczą one prawie 40\% globalnych zbiorów żywności\cite{FAO2025PlantDiseasesFoodSecurity}. Przekłada się to na wielomiliardowe straty. Poza aspektem ekonomicznym, warto zwrócić uwagę na aspekt społeczny - w samej Polsce, stan na 2025 rok, nawet 1,4 miliona mieszkańców zmaga się z głodem \cite{PCK2025Glod}. Straty spowodowane chorobami zwiększają ceny żywności w sklepach, pogłębiając problem dostępu do jedzenia, szczególnie dla osób ubogich. W związku z tym ważne jest opracowanie i zbadanie rozwiązań detekcji chorób roślin, w celu zmniejszenia strat oraz zapobiegania masowemu rozprzestrzenianiu się chorób.
\subsection{Cel i zakres pracy}
Celem tej pracy jest zbadanie, jak wybrane architektury głębokiego uczenia różnią się pod względem skuteczności oraz odporności na redukcję liczby danych treningowych i degradację ich jakości w zadaniu klasyfikacji chorób roślin. Dodatkowym celem jest stworzenie aplikacji mobilnej, stanowiącej przykład praktycznej implementacji badanej technologii. W ramach pracy  opracowano badania umożliwiające ocenę dostępnych publicznie modeli. Przygotowano także demo technologiczne.
\subsection{Układ pracy} 
Praca zawiera dziesięć rozdziałów. Rozdział pierwszy to wstęp zawierający sformowanie problemu, cel i zakres pracy.

Rozdział drugi przedstawia przegląd literatury, prezentując badania w temacie, zastosowanie sztucznej inteligencji w rolnictwie, zbiory danych zawierające zdjęcia oraz przykłady aplikacji, które służą do detekcji chorób roślin.

Rozdział trzeci jest wprowadzeniem do uczenia maszynowego. Przedstawia podstawowe pojęcia potrzebne do zrozumienia pracy oraz przegląd technik uczenia maszynowego.

Rozdział czwarty poświęcony jest opisowi wybranych architektur poddanych analizie. 

Rozdział piąty przedstawia wybrany zbiór danych do badań.

W ramach rozdziału szóstego przedstawiono opracowane badanie. Opisano użyte techniki oraz scenariusze eksperymentów. 

Rozdział siódmy przedstawia wyniki eksperymentów.

Rozdział ósmy poświęcony jest analizie uzyskanych wyników.

Rozdział dziewiąty przedstawia wykonane demo technologiczne oraz wykorzystane do tego technologie.

Ostatni, dziesiąty rozdział stanowi podsumowanie całej pracy oraz przedstawia możliwe kierunki rozwoju badań.
\clearpage